% ============================================================================
% Annexe D — Glossaire technique
% ============================================================================
\chapter{Glossaire technique}

\begin{description}[leftmargin=3.5cm, style=nextline, font=\bfseries\color{NavyBlue}]

  \item[AIOps]
    \textit{Artificial Intelligence for IT Operations}. Application des techniques
    d'apprentissage automatique et d'IA aux opérations informatiques, notamment
    pour la détection d'anomalies, la corrélation d'événements et la génération
    automatique de recommandations.

  \item[Alert fatigue]
    (\textit{fatigue d'alertes}) Phénomène psychologique par lequel les opérateurs
    finissent par ignorer les alertes en raison de leur volume excessif ou d'un
    taux élevé de faux positifs.

  \item[Asyncio]
    Bibliothèque Python de programmation asynchrone basée sur des coroutines
    (async/await) et une boucle d'événements (\textit{event loop}), permettant la
    gestion concurrente de nombreuses connexions réseau sans threads système.

  \item[Chunk]
    Dans TimescaleDB, unité de stockage physique correspondant à une partition
    temporelle de l'hypertable. Les chunks anciens sont automatiquement compressés.

  \item[Embedding]
    Représentation vectorielle dense d'un texte dans un espace de haute dimension.
    Deux textes sémantiquement proches produisent des vecteurs proches selon la
    distance cosinus. Utilisé dans \caimp{} pour la recherche de runbooks et
    d'incidents similaires.

  \item[GGUF]
    Format de fichier pour les modèles de langage quantifiés, utilisé par Ollama
    et \texttt{llama.cpp}. Remplace les formats GGML et GGML-v3.

  \item[Goroutine]
    Fil d'exécution léger du langage Go, géré par le runtime Go plutôt que par
    l'OS. Des milliers de goroutines peuvent s'exécuter simultanément avec
    une empreinte mémoire minimale.

  \item[Healthcheck]
    Dans Docker, commande exécutée périodiquement pour vérifier que le service
    est opérationnel. Si le healthcheck échoue plusieurs fois, Docker peut
    redémarrer le conteneur.

  \item[HNSW]
    \textit{Hierarchical Navigable Small World}. Algorithme de recherche approximative
    de voisins les plus proches dans les espaces vectoriels de haute dimension.
    Offre une complexité de recherche $O(\log n)$ avec une précision proche
    de la recherche exhaustive.

  \item[Holt-Winters]
    Algorithme de lissage exponentiel double (niveau + tendance) pour la prévision
    de séries temporelles à tendance linéaire. Conçu par Charles Holt (1957) et
    Peter Winters (1960).

  \item[Hypertable]
    Extension de table PostgreSQL par TimescaleDB qui partitionne automatiquement
    les données par intervalles de temps, permettant des performances d'insertion
    et de requête très élevées sur les données temporelles.

  \item[JWT]
    \textit{JSON Web Token}. Standard RFC 7519 pour la représentation de claims
    entre deux parties sous forme de JSON signé (JWS) ou chiffré (JWE). Utilisé
    dans \caimp{} pour l'authentification des requêtes API.

  \item[LLM]
    \textit{Large Language Model}. Modèle de réseau de neurones entraîné sur de
    grandes quantités de texte pour générer ou compléter du texte en langage
    naturel. Exemples : Llama, GPT, Mistral, Gemma.

  \item[mTLS]
    \textit{mutual Transport Layer Security}. Extension du protocole TLS dans
    laquelle les deux parties (client et serveur) s'authentifient mutuellement
    par certificat numérique. Utilisé dans \caimp{} pour authentifier les agents
    de supervision.

  \item[NATS JetStream]
    Extension persistante du système de messagerie NATS, offrant des garanties
    de livraison at-least-once, des consommateurs durables et un stockage
    des messages sur disque.

  \item[Ollama]
    Serveur d'inférence LLM local exposant une API REST. Prend en charge les
    modèles GGUF (quantifiés) et permet l'exécution de modèles sur des machines
    sans GPU.

  \item[OTLP]
    \textit{OpenTelemetry Protocol}. Protocole standardisé par la CNCF pour la
    transmission de métriques, traces et journaux entre les agents d'instrumentation
    et les backends d'observabilité.

  \item[pgvector]
    Extension PostgreSQL ajoutant le type de données \texttt{VECTOR(n)} et les
    opérateurs de similarité (cosinus, produit scalaire, distance L2). Permet
    la recherche de voisins les plus proches directement dans PostgreSQL.

  \item[Q4\_K\_M]
    Schéma de quantification 4 bits utilisant la méthode K-means pour grouper
    les poids du modèle. Offre un excellent compromis entre taille (réduction
    de 75~\% vs fp16) et qualité d'inférence.

  \item[RAG]
    \textit{Retrieval-Augmented Generation}. Technique qui améliore la qualité
    des réponses d'un LLM en récupérant dynamiquement des documents pertinents
    depuis une base de connaissances et en les injectant dans le prompt.

  \item[RLS]
    \textit{Row-Level Security}. Fonctionnalité PostgreSQL permettant de définir
    des politiques de contrôle d'accès au niveau de chaque ligne d'une table.
    Utilisée dans \caimp{} pour l'isolation multi-tenant.

  \item[SSE]
    \textit{Server-Sent Events}. Protocole HTTP unidirectionnel permettant
    au serveur d'envoyer des flux de données en continu vers le client.
    Utilisé dans \caimp{} pour le streaming des réponses du chat IA.

  \item[Subject NATS]
    Chaîne de caractères hiérarchique identifiant le canal d'un message NATS
    (ex: \texttt{anomaly.detected.org-1.critical}). Les abonnements peuvent
    utiliser des jokers (\texttt{>} pour zéro ou plusieurs niveaux).

  \item[Time to critical]
    (\textit{délai jusqu'à la criticité}) Dans \caimp{}, nombre d'heures prévu
    avant qu'une métrique dépasse le seuil de 90~\%, calculé par interpolation
    sur les prévisions Holt-Winters.

  \item[TimescaleDB]
    Extension PostgreSQL transformant une base de données relationnelle en base
    de données temporelles haute performance, avec partitionnement automatique
    (hypertables), compression columnar et politiques de rétention automatiques.

  \item[VPS]
    \textit{Virtual Private Server}. Serveur virtuel isolé hébergé dans un
    datacenter tiers, proposé par des fournisseurs comme OVH, DigitalOcean
    ou Hetzner.

  \item[WebSocket]
    Protocole de communication bidirectionnel full-duplex sur une connexion TCP
    persistante (RFC 6455). Utilisé dans \caimp{} pour pousser les événements
    en temps réel vers les navigateurs connectés.

  \item[Z-score]
    Mesure statistique exprimant l'écart d'une observation à la moyenne en
    unités d'écart-type :
    $z = (x - \mu) / \sigma$. Une valeur $|z| > 3$ est considérée comme
    statistiquement anormale dans une distribution normale.

\end{description}
